\chapter{Журнал и инварианты}

Журнал есть плоскость управления в чистом виде: всё, что имеет силу в системе, существует как запись в нём. Эта глава описывает форму журнала, а затем разбирает десять инвариантов, потому что именно журнал делает их проверяемыми.

\section{Запись}

Запись содержит хеш предыдущей записи, тип, тело или хеши тел, и подписи. Тело сериализуется в канонический JSON: детерминированная форма, в которой одинаковые данные всегда дают одинаковые байты, а значит одинаковый хеш. Подписи Ed25519 ставятся поверх канонической формы; идентификатор ключа есть первые шестнадцать шестнадцатеричных символов хеша самого ключа. Набор пинов кор-ключей ротируется через перекрытие: сначала распространяется набор со старым и новым ключом, потом подписи переходят на новый, потом старый удаляется. Скомпрометированный ключ удаляется немедленно.

\begin{exbox}[{Запись назначения, набросок}]
\begin{lstlisting}
{
  "type": "assignment",
  "prev": "b3c1...",
  "task": "TASK-260901-7fk2",
  "node": "ed25519:...",
  "role": "producer",
  "plan_digest": "9a4e...",
  "lease": {"epoch": 41, "token": "f81d..."},
  "slot": "env-1/slot-07",
  "body_hash": "51aa...",
  "availability": ["replica-a", "replica-b"],
  "sig": [{"key_id": "0123456789abcdef", "signature": "..."}]
}
\end{lstlisting}
\end{exbox}

\section{Цепочка, дерево и снапшот}

Для записи $i$ с канонической формой $c(r_i)$ хеш вычисляется как $h_i = H(h_{i-1} \,\|\, c(r_i))$; у первой записи предшественник есть строка нулей. Над хешами записей строится дерево Меркла, и снапшот, подписанный кор-ключами, содержит корень дерева, длину журнала, голову и версию. История только дописывается: удаление и замена запрещены.

Что это даёт практически. Любой участник с последним снапшотом проверяет любую запись по доказательству включения. Изменение прошлого невозможно без пересчёта всех последующих хешей и новой подписи снапшота, то есть без кворума. Светлые клиенты, среда и контрибьюторы, не хранят журнал целиком.

\section{Типы записей}

\begin{center}
\small
\begin{tabular}{p{3.4cm} p{10.2cm}}
\toprule
Тип & Что фиксирует \\
\midrule
цель & область изменений, класс риска, бюджет; подпись кор-ноды \\
capability & часть цели, выданная оркестратору: пути, типы изменений, лимиты \\
контракт задачи & версия: write scope, зависимости, критерии, профиль свойств, оценка, класс, бюджеты; digest плана \\
зависимость & какие принятые артефакты задача использует и какие результаты предоставляет \\
назначение & задача, нода, роль, аренда, слот \\
ревизия & хеш коммита, база, голова цепочки receipts, подпись ноды \\
заморозка кандидатов & множество нод, из которых будет выбран ревьювер \\
случайность & значение маяка или VRF после заморозки и ревизии \\
кандидат & база $M$, неизменяемый кандидат $I$ \\
ревью кандидата & одобрение или находки, адресованные к $I$, версии контракта, политике, окружению \\
аттестация & свойство из профиля, исход (пройдена, провалена, не установлена), адрес, нода, верификатор при наличии \\
интеграция & пара $(M, I)$, ссылки на ревью и аттестацию \\
вестинг & закрепление начисления после $k$ эпох \\
аренда, эпоха & нода, эпоха, токен, предшественник, причина, чекпоинт, область \\
чекпоинт & хеш снимка сессии с сертификатом доступности \\
пауза, возобновление & экстренная остановка и снятие \\
состав & изменение множества кор-нод и ключей \\
операция раннера & авторизация именованной операции с идентификатором и ключом идемпотентности; отдельной записью исход \\
аттестация репутации & исход задачи для ноды \\
подтверждение идентичности & идентичность, класс доказательства, хеш свидетельства резолвера, окно; подписи валидаторов, каждый резолвил сам \\
сертификация ключа & ключ ноды, подписанный ключом идентичности, с началом и концом действия \\
\bottomrule
\end{tabular}
\end{center}

\section{Порядок, кворум и уровни}

Порядок записей задаёт секвенсер. Запись совершена, когда собран сертификат кворума уровня. На L0 секвенсер и единственный валидатор совпадают. На L1 секвенсер есть всегда включённый сервис проекта, которому доверяют в отношении порядка, а кор-ноды подписывают асинхронно; запись с $t$ подписями, $t > n/2$, совершена, и до этого момента среда её не исполняет. Порог большинства защищает от одиночного несанкционированного решения, но не от раздвоения истории при злонамеренном секвенсере (глава 3); L1 это признаёт, а не скрывает. На L2 порядок фиксируется византийски устойчивым консенсусом валидаторов, который даёт и то и другое при $f$ лжецах; среда исполняет только записи с $2f+1$ подписями.

\begin{whybox}[Почему секвенсер отдельно от валидаторов]
Консенсус между ноутбуками кор-команды непрактичен: живучесть требует, чтобы достаточное число нод было в сети постоянно. Секвенсер даёт живучесть, валидаторы дают безопасность. Так устроены роллапы, и так уже устроен реестр аудита, на котором журнал будет реализован.
\end{whybox}

Роль лидера консенсуса на L2 механическая: она ротируется по раундам, не требует модели и не совпадает с ролью оркестратора. Это важно повторять: оркестратор не лидер журнала, он один из отправителей записей.

\section{Доступность и материализация}

Тела задач, находки ревью, снимки сессий и артефакты хранятся адресуемо по содержимому, журнал содержит только хеши. Хеш включается в журнал только с сертификатом доступности: несколько независимых реплик подтвердили хранение объекта. Код тоже адресуется по содержимому: журнал хранит хеши коммитов, а git выступает хранилищем, в котором любая нода проверяет, что ревью, тесты и интеграция относятся к одному состоянию.

Журнал может материализоваться как борд задач в git-репозитории проекта: каждый переход есть коммит борда, сертификат кворума покрывает хеш коммита, полная история задачи есть история git. Журнал не дублирует борд, он сертифицирует его. Для команды это означает, что привычный борд остаётся и что ничего не нужно переносить: добавляются подписи и цепочка. То же относится к графу зависимостей задач и артефактов: он материализуется из записей о зависимостях и не является вторым источником истины.

\section{Десять инвариантов}

Каждый инвариант ниже сопровождён тем, что ломается без него. Это полезный тест при реализации: если предлагаемое упрощение делает поломку возможной, оно неприемлемо.

\begin{invbox}[I1. Эффект только по записи]
Ни одного эффекта на среде без совершённой записи авторизации, digest плана которой совпадает с digest плана вызова, и актуальность которой проверена на границе исполнения.

Без него: любой, кто дотянулся до среды, исполняет что угодно; оркестратор в одиночку открывает санбоксы; агент с расширенным планом работает под чужой задачей; операция, разрешённая в прошлой эпохе, исполняется после остановки.
\end{invbox}

\begin{invbox}[I2. Только сужение]
План задачи есть подмножество capability, capability есть часть цели; ревизия не изменяет ничего вне write scope задачи; условия приёмки меняются только новой версией контракта с отдельной авторизацией.

Без него: оркестратор создаёт задачи с полномочиями, которых цель не давала; отравленный субагент расширяет права родителя; задача трогает тесты, которые должна была только читать; критерий успеха ослабляется по ходу работы.
\end{invbox}

\begin{invbox}[I3. Интегрируется ровно проверенное]
Основная ветка продвигается только атомарной заменой $M$ на неизменяемого кандидата $I$, к которому адресованы финальное одобрение и аттестации, ссылающиеся на текущую версию контракта и политики.

Без него: две безопасные по отдельности задачи сливаются в небезопасную композицию, которую никто не проверял; одобрение, сделанное на старой базе, переезжает на кандидата над новой головой. Исполняемая модель в главе 11 показывает оба сценария как нарушения.
\end{invbox}

\begin{invbox}[I4. Суждение вне консенсуса]
Валидаторы проверяют только детерминированные предикаты над записями и хешами.

Без него: консенсус о правильности кода невозможен и превращается в политику; византийская устойчивость теряет смысл, потому что предмет голосования не верифицируем.
\end{invbox}

\begin{invbox}[I5. Независимость и случайность]
Продьюсер, ревьювер и исполнитель проверок независимы согласно объявленной политике проекта; отсутствие подходящих участников не ослабляет политику автоматически; случайность выбора фиксируется после заморозки кандидатов и хеша кандидата.

Без него: сговор продьюсера и ревьювера; одинаковые слепые зоны одной модели; подбор ревьювера оркестратором; тихое ослабление требований, когда некому ревьюить.
\end{invbox}

\begin{invbox}[I6. Ключи у владельца, секреты в раннере]
Ключи контрибьютора не покидают его машину, и подпись выполняется на машине агентского цикла; секреты проекта существуют только в раннере и недоступны ни из одного санбокса.

Без него: компрометация среды означает компрометацию всех контрибьюторов; ключи подписи проекта извлекаются через shell.
\end{invbox}

\begin{invbox}[{I7. Журнал истина, сессия кэш}]
Записи с устаревшей арендой отвергаются и сохраняются как расхождение.

Без него: два оркестратора одновременно; отравленный контекст переезжает вместе с ролью; смена оркестратора теряет состояние.
\end{invbox}

\begin{invbox}[{I8. Остановка дешёвая, откат невозможен}]
Любая кор-нода останавливает, кворум возобновляет; среда хранит наивысшую принятую версию журнала.

Без него: скомпрометированный оркестратор блокирует остановку; остановку обходят старым снапшотом.
\end{invbox}

\begin{invbox}[I9. Ограничение как факт]
Среда объявляет, что способна ограничить; задача, требующая полного ограничения, отказывает на частичном.

Без него: задача, рассчитанная на microVM, исполняется под Seatbelt, и никто об этом не знает.
\end{invbox}

\begin{invbox}[I10. Происхождение обращений]
Все обращения к среде через протокол и выданные ею результаты имеют проверяемое происхождение и связаны цепочкой receipts, голова которой входит в запись о ревизии. Протокол не удостоверяет полноту внутреннего контекста недоверенной рабочей ноды.

Без него: нельзя установить, что именно среда исполнила по чьему запросу; нельзя сравнить результат с верификатором. Вторая фраза не ослабление, а честная граница: во второй редакции инвариант обещал подотчётность всего, что модель видела, чего протокол наблюдать не может.
\end{invbox}
